%%%%%%%%%%%%%%%%%%%%%%%%%%%%%%%%%%%%%%%%%%%%%%%%%%%%%%%%%%%%%%%%%%%%%%
% Lớp 11 | Đường thẳng, mặt phẳng. Quan hệ song song trong | Hai đường thẳng song song
% Dạng 1H4?2-8: Toán thực tế áp dụng hai đường thẳng song song
%%%%%%%%%%%%%%%%%%%%%%%%%%%%%%%%%%%%%%%%%%%%%%%%%%%%%%%%%%%%%%%%%%%%%%


% nguồn: 2026-11-15_11_1H4_TL2_Song_song (nhập 15/11/2026)
\begin{bt}%[1H4H2-8]%[ma:314451fc]
% nguon: Archive/2024/TOAN_11_XE/TOAN_11_XE/data4/11C4/KNTT-4.2.tex #779

	Quan sát phòng học của lớp và nêu lên hình ảnh của hai đường thẳng song song, cắt nhau, chéo nhau.
	\loigiai{Tùy theo tình hình thực tế lại lớp, chẳng hạn như hai mép cửa chính, lề bảng và mép tường,...}
\end{bt}

% nguồn: 2026-11-15_11_1H4_TL2_Song_song (nhập 15/11/2026)
\begin{bt}%[1H4H2-8]%[ma:6cb3cf59]
% nguon: Archive/SP_KhoLatex/ID6_old/1-GK1-CanhDieu-De6-NH23-24.tex #3026

	\immini{
			Một bể kính chứa nước có đáy là hình chữ nhật được đặt nghiêng như hình vẽ. Giải thích tại sao đường mép nước $ E F $ song song với cạnh $ C D $ của bể nước?
	}{
			\begin{tikzpicture}[scale=0.7, font=\footnotesize,>=stealth]%<DTools>
			\def\canhAD{4};\def\canhBA{2};\def\gocBAD{-130};\def\h{4};\def\xdinhA'{1};
			%Định nghĩa điểm.
			\coordinate (A) at (0,0);
			\coordinate (B) at ($(A)+(\gocBAD:\canhBA)$);
			\coordinate (C) at ($(B)+(0:\canhAD)-(0,1)$);
			\coordinate (D) at ($(A)+(0:\canhAD)-(0,1)$);
			\coordinate (A') at ($(A)+(\xdinhA',\h)$);
			\coordinate (B') at ($(B)+(\xdinhA',\h)$);
			\coordinate (C') at ($(C)+(\xdinhA',\h)$);
			\coordinate (D') at ($(D)+(\xdinhA',\h)$);
			\coordinate (E) at ($(D)!1/3!(D')$);
			\coordinate (F) at ($(C)!1/3!(C')$);
			\fill[cyan!50, opacity=0.5] (A)--(B)--(F)--(E)--cycle;
			\fill[cyan, opacity=0.5] (A)--(B)--(C)--(D)--cycle;
			\fill[cyan, opacity=0.5] (E)--(F)--(C)--(D)--cycle;
			%Vẽ khối lẳng trụ ABCD.A'B'C'D'.
			\draw (A')--(B')--(B)--(C)--(C')--(D')--cycle (B')--(C') (D')--(D)--(C) (B)--(F)--(E);
			\draw[dashed] (A)--(D) (A')--(A)--(B) (A)--(E);
			%Hiển thị các điểm.
			\foreach \x/\y in {C/0, D/0, E/0, F/0}{\fill (\x) circle(1pt) ($(\x)+(\y:0.3cm)$) node{$\x$};}
		\end{tikzpicture}
	}

	\loigiai{
		
		\immini{
			Giả sử mặt phẳng $ (A B E F) $ là mặt nước, mặt phẳng $ (A B C D) $ là mặt đáy của bể kính và $ (E F C D) $ là mặt bên của bể kính. Ba mặt phẳng $ (A B E F), (E F C D), (A B C D) $ là ba mặt phẳng cắt nhau theo 3 giao tuyến $ E F, A B, C D $. Vì $ C D \parallel E F $ (do đáy bể là hình chữ nhật) nên ba đường thẳng $ E F, A B, C D $ đôi một song song.\\
			Vậy đường mép nước $ EF $ song song với cạnh $ C D $ của bể nước.
		}{
			\begin{tikzpicture}[scale=0.7, font=\footnotesize,>=stealth]%<DTools>
				\def\canhAD{4};\def\canhBA{2};\def\gocBAD{-130};\def\h{4};\def\xdinhA'{1};
				%Định nghĩa điểm.
				\coordinate (A) at (0,0);
				\coordinate (B) at ($(A)+(\gocBAD:\canhBA)$);
				\coordinate (C) at ($(B)+(0:\canhAD)-(0,1)$);
				\coordinate (D) at ($(A)+(0:\canhAD)-(0,1)$);
				\coordinate (A') at ($(A)+(\xdinhA',\h)$);
				\coordinate (B') at ($(B)+(\xdinhA',\h)$);
				\coordinate (C') at ($(C)+(\xdinhA',\h)$);
				\coordinate (D') at ($(D)+(\xdinhA',\h)$);
				\coordinate (E) at ($(D)!1/3!(D')$);
				\coordinate (F) at ($(C)!1/3!(C')$);
				\fill[cyan!50, opacity=0.5] (A)--(B)--(F)--(E)--cycle;
				\fill[cyan, opacity=0.5] (A)--(B)--(C)--(D)--cycle;
				\fill[cyan, opacity=0.5] (E)--(F)--(C)--(D)--cycle;
				%Vẽ khối lẳng trụ ABCD.A'B'C'D'.
				\draw (A')--(B')--(B)--(C)--(C')--(D')--cycle (B')--(C') (D')--(D)--(C) (B)--(F)--(E);
				\draw[dashed] (A)--(D) (A')--(A)--(B) (A)--(E);
				%Hiển thị các điểm.
				\foreach \x/\y in {A/180, B/180, C/0, D/0,E/0, F/0}{\fill (\x) circle(1pt) ($(\x)+(\y:0.3cm)$) node{$\x$};}
			\end{tikzpicture}
		}
		 } 
\end{bt}
